\section{Data and measurement}
\subsection{Survey}

The recorded survey interviews span 9 December 2025 to 17 February 2026. Fundación Capital recruited migrant women leaders, who reached potential participants through contacts and referrals and administered the survey. Recruitment was not restricted to former program participants. The intended population included adult women with international migration experience, including Colombian return migrants, with about three months' residence after arrival or return. The programmed screen used consent, gender, age and receipt of remittances by the respondent or someone in her household in the previous 12 months. The form did not test every intended recruitment condition. Household receipt does not establish that the respondent personally received or managed a transfer.

The revision retains the historical 423-record cohort and the original manual and coded exclusions. This is a nonprobability sample. Referral and household links were not systematically recorded, the cohort is not a count of verified distinct households, and there is no denominator for a response rate. Residence is recorded residence, not the interview site. The proposed cohort table (T1) describes the admitted records and the form-version partition; any additional geography panel has a separate source and disclosure check. A screening-flag diagnostic does not define a second set of exclusions.

The first outcome distinguishes respondents reporting any digital retention or use of remittance proceeds from those reporting withdrawal of all proceeds in cash. We also retain the five response categories in questionnaire order. Two choices describe amounts on the same side of the halfway point in different words. We do not reconstruct monetary shares or treat the choices as a strictly ranked scale. The latest receipt channel and the 12-month receipt frequency are separate context items; the latter preserves the combined response for nonapplicability or not remembering.

The second outcome is whether respondents report making or receiving a domestic digital payment or transfer in the previous 60 days. The denominator is valid yes/no answers, with refusal and other unscored states kept separate. Among affirmative domestic users, we retain the original five operation bins on the valid-bin base and report the response-state distribution on the affirmative-user base. Uncertainty is not a bin. Shared endpoints in the response labels cannot be resolved from the data. Fee visibility concerns the last operation in a separate 90-day window. The proposed T2 and F1 display the two outcomes without a common denominator or payment-flow arrows.

The working models use three survey measures. The access/basic-use measure combines reported phone access/use, internet-use frequency and running out of data or balance; some content is close to the outcomes. The second averages reported ability to send or receive money by phone and confidence in avoiding online fraud and suspicious messages. It measures neither tested skills nor confidence in making a transfer as a standalone construct. The third combines address documentation in the respondent's name and a phone line in her own name; these are two onboarding conditions, not account ownership or provider-verified eligibility. All specified components must be present. The optional three-item alternative adds app installation under its own applicability rule. Pair-complete and common-complete comparisons are distinct; Pearson correlations do not validate a scale.

\subsection{Interviews, focus groups and documentary comparison}

The original methods chapter reports 30 semistructured interviews with Venezuelan migrant women and six focus groups involving 47 participants, across Bogotá, Soacha, Cali and Medellín. These are reported activity counts, not a claim about the number of distinct people across activities. The qualitative analysis combined themes informed by the study's conceptual domains with themes developed from the interview and focus-group material. We retain the existing authored findings. We do not infer their frequency in the survey or assume participant-level linkage.

The original paper also draws on institutional and regional expert interviews and reviews regulatory, central-bank and policy documents. Experts describe implementation and institutional responsibilities; official sources document selected rules and system designs. Neither is evidence that a particular respondent could open an account, used a specific rail or resolved a complaint. The regional comparison examines counterpart systems, not comparable groups of women in five countries.

\exhibitplaceholder{T1}{Cohort and form versions}{Admitted total and version accounting. Screening states are a separate diagnostic, not exclusions. No geography enters from the cohort emitter alone.}
